On a Poisson manifold, two points are called **equivalent** if they can be joined by a path consisting of segments of integral curves of Hamiltonian vector fields. The equivalence classes under this relation are called the **leaves** (or **symplectic leaves**) of the Poisson manifold.

Each leaf is an immersed submanifold carrying a natural symplectic structure, and the Poisson bracket of the ambient manifold restricts to the Poisson bracket induced by this symplectic structure on each leaf. The values of all possible Casimir functions (functions whose Poisson bracket with everything vanishes) are constant on each leaf.

The leaves form a (generally singular) foliation of the Poisson manifold. In general, they are not a regular foliation: leaves may have different dimensions. The dimension of a leaf through a point equals the rank of the Poisson bivector field at that point.

**Example.** For the Poisson bracket of angular momentum components $\{F, G\} = \sum (\partial F/\partial M_i)(\partial G/\partial M_j)\{M_i, M_j\}$ on $\mathbb{R}^3$, the leaves are the concentric spheres $M_1^2 + M_2^2 + M_3^2 = \text{const}$ (symplectic manifolds of dimension 2) and the origin (a single point, dimension 0).
